\documentclass[11pt,a5paper]{article}
\usepackage[margin=16mm]{geometry}
\usepackage{amsmath,amssymb,amsthm,lmodern}
\newtheorem{theorem}{Theorem}
\title{A note on rank-one structure}
\author{}\date{}
\begin{document}
\begin{center}{\Large\bfseries A note on rank-one structure}\end{center}
\section{One direction}
A rank-one matrix has a simple geometry.
For $A=uv^{\mathsf T}$, every output lies along $u$;
only the component of the input along $v$ matters.

\begin{theorem}\label{thm:rank-one}
For $u,v\in\mathbb R^d$, the operator norm satisfies
\[\|uv^{\mathsf T}\|_{\mathrm{op}}=\|u\|_2\,\|v\|_2.\]
\end{theorem}
\begin{proof}
For every unit vector $x$, Cauchy--Schwarz gives
\[\|Ax\|_2=\|u\|_2\,|v^{\mathsf T}x|
  \leq\|u\|_2\,\|v\|_2.\]
When $v\ne0$, take $x=v/\|v\|_2$ to attain equality.
The case $v=0$ is immediate.
\end{proof}

\section{Beyond rank one}
Theorem~\ref{thm:rank-one} identifies the scale of a single component.
For a sum of components, the triangle inequality yields
\[\Bigl\|\sum_{j=1}^r u_jv_j^{\mathsf T}\Bigr\|_{\mathrm{op}}
 \leq\sum_{j=1}^r\|u_j\|_2\,\|v_j\|_2.\]
No orthogonality is assumed. Sharper conclusions
depend on how the directions interact.
\end{document}
